% ============================================================
%   CHAPTER 4: Description of the Design Elements
% ============================================================
\chapter{Description of the Design Elements}
\label{chap:design_elements}

\begin{table}[H]
    \caption{Revision History – Chapter 4}
    \centering
    \begin{tabular}{llll}
        \toprule
        Version & Date             & Author     & Description                            \\
        \midrule
        1.0     & January 22, 2026 & Y. Sangale & Initial Draft                          \\
        \midrule
        1.1     & April 21, 2026   & Y. Sangale & Added FPU design element (Section 4.2) \\
        \bottomrule
    \end{tabular}
\end{table}

\textcolor{blue}{\textbf{\textit{This Section will be completed later.}}}

This chapter provides a detailed technical description of all major design
elements composing the \gls{rv64i} 64-bit single-core microcontroller. Each
sub-block is presented in terms of hardware interfaces, bus connectivity,
internal registers, and its functional and structural roles within the
\gls{SoC}. All diagrams and tables are placeholders for later detailed
schematics.

% ============================================================
\section{\gls{cpu} Core}

\subsection{HW Interfaces to Other Blocks}
The \gls{cpu} Core is the central processing unit of the system, implementing
the RV64IM instruction set and acting as the single bus master in the
microcontroller architecture. It connects to all major subsystems as follows:
\begin{itemize}
    \item \textbf{Instruction Memory (\gls{I-MEM}):} Connected directly via the \textbf{Instruction Bus (I-Bus)}
          for program fetch under the \gls{Harvard} architecture model.
    \item \textbf{Data Memory (\gls{D-MEM}):} Connected directly via the \textbf{Data Bus (D-Bus)} for load/store access.
    \item \textbf{Wishbone Bus:} Acts as the \textbf{bus master} interfacing to all memory-mapped peripherals
          (\gls{gpio}, \gls{uart}, \gls{qspi}, Interrupt Controller, and Clock Unit).
    \item \textbf{Interrupt Input:} Receives a single aggregated interrupt signal from the Vectored Interrupt Controller (VIC).
    \item \textbf{Clock and Reset:} Driven by the system clock from the Clock Generation Unit (CGU) and
          synchronized reset signal.
\end{itemize}

\begin{figure}[H]
    \centering
    \caption{\gls{cpu} Core interface diagram showing I-Bus, D-Bus, and Wishbone connections to system blocks.}
    \label{fig:cpu_interface}
\end{figure}

\subsection{Bus Interface}
The \gls{cpu} operates as a \textbf{Wishbone Master} for all peripheral
transactions. It initiates read/write cycles with standard Wishbone control
signals: \texttt{ADR\_O}, \texttt{DAT\_O}, \texttt{DAT\_I}, \texttt{STB\_O},
\texttt{ACK\_I}, \texttt{CYC\_O}, and \texttt{WE\_O}. Instruction and data
memories are accessed directly via the independent I-Bus and D-Bus, each 64-bit
wide for full datapath utilization.

\begin{figure}[H]
    \centering
    \caption{\gls{cpu} bus architecture — direct \gls{Harvard}-style connections to \gls{I-MEM}/\gls{D-MEM} and Wishbone master for peripherals.}
    \label{fig:cpu_bus_arch}
\end{figure}

\subsection{Registers}
The \gls{cpu} core implements 32 × 64-bit general-purpose registers
(\texttt{x0–x31}) and a 64-bit program counter. Several internal control and
status registers are included for interrupt and \gls{pipeline} management.

\begin{table}[H]
    \centering
    \caption{\gls{cpu} Core Register Summary}
    \label{tab:cpu_registers}
    \begin{tabular}{|l|l|}
        \hline
        \textbf{Register} & \textbf{Description}                                     \\ \hline
        \gls{pc}          & 64-bit Program Counter for sequential and branch control \\ \hline
        IR                & Instruction Register holding current fetched instruction \\ \hline
        SR                & Status Register with control and flag bits               \\ \hline
        IBASE             & Base address for interrupt vector table                  \\ \hline
        MSTATUS           & Machine Status Register (\gls{rv64i} subset)             \\ \hline
    \end{tabular}
\end{table}

\subsection{Functional Overview}
The \gls{cpu} executes instructions using a deterministic scalar \gls{pipeline}
that supports the \gls{riscv} RV64IM base instruction set. Internally, it
consists of the following major sub-modules:
\begin{itemize}
    \item \textbf{Instruction Decoder:} Decodes opcodes fetched from \gls{I-MEM} and generates control signals for \gls{alu} and memory.
    \item \textbf{Program Counter (\gls{pc} ):} Maintains the address of the current instruction and computes sequential or branch targets.
    \item \textbf{\gls{alu}:} Performs arithmetic, logical, shift, and comparison operations; supports integer multiply/divide.
    \item \textbf{Register File:} Provides two read ports and one write port for operand access and result storage.
    \item \textbf{Interrupt Input:} Interfaces with VIC to handle asynchronous interrupt requests.
\end{itemize}

Each instruction is executed in a simple in-order flow without speculation or
\gls{pipeline} hazards beyond single-cycle stalls.

\begin{figure}[H]
    \centering
    \caption{\gls{cpu} Core internal structure showing Instruction Decoder, \gls{alu}, Program Counter, and Register File.}
    \label{fig:cpu_structure}
\end{figure}

\subsection{Structural Overview}
The \gls{cpu} architecture follows a \gls{Harvard} model with:
\begin{itemize}
    \item Dedicated instruction path (I-Bus) for fetching 64-bit instructions.
    \item Dedicated data path (D-Bus) for memory and peripheral data transactions.
    \item Unified Wishbone interface for peripheral control and configuration.
\end{itemize}

The internal control path coordinates fetch, decode, execute, and memory
operations. Instruction flow synchronization ensures deterministic timing
suitable for embedded applications.

\begin{figure}[H]
    \centering
    \caption{Simplified \gls{Harvard} architecture connection of \gls{cpu} to \gls{I-MEM}, \gls{D-MEM}, and Wishbone peripherals.}
    \label{fig:cpu_harvard_arch}
\end{figure}

\subsection{Interrupts}
A single interrupt input line connects to the VIC. The \gls{cpu} jumps to a
vector address computed as:
\[
    \text{ISR Vector} = \text{IBASE} + (\text{\gls{IRQ}\_ID} \times 0x10)
\]
Upon completion, control returns to the instruction following the interrupted
instruction.

\subsection{Clock and Reset}
The \gls{cpu} operates at an 80\,MHz clock supplied by the CGU. Reset
initializes the \gls{pc} to 0x0000\_0000, clears all \gls{pipeline} stages, and
resets the internal status registers.

\subsection{Voltage Class}
Core operates in the 3.3\,V logic domain.

\section{Floating-Point Unit (FPU)}

\subsection{Introduction}
This section specifies the Floating-Point Unit (FPU) added to the RV64I\_RWU microcontroller as an integral extension. The FPU implements the RISC-V D-extension (RV64D) and is fully compliant with the IEEE 754-2008 Standard for Floating-Point Arithmetic in double-precision (64-bit) format. The FPU is tightly coupled to the CPU pipeline, sharing the instruction fetch and decode stages, and operates on a dedicated 64-bit floating-point register file (f0–f31) separate from the integer register file.

The FPU is designed for deterministic, in-order execution consistent with the single-pipeline philosophy of the microcontroller. All floating-point operations are performed in hardware with no software emulation fallback. The FPU supports all mandatory IEEE 754-2008 operations, all five rounding modes, and all five floating-point exception flags.

\subsection{Compliance and Standards}
The FPU is designed to comply with the following standards and specifications:

\begin{table}[H]
\centering
\begin{tabular}{|l|l|l|}
\hline
\textbf{Standard / Specification} & \textbf{Version / Reference} & \textbf{Compliance Level} \\ \hline
IEEE 754-2008 & Standard for Floating-Point Arithmetic & Full (double-precision) \\ \hline
RISC-V Unprivileged ISA & Volume I, D-Extension (RV64D) & Full \\ \hline
RISC-V Unprivileged ISA & Volume I, F-Extension (RV64F) & Full (subset of D) \\ \hline
RISC-V Privileged ISA & fcsr register (frm + fflags) & Full \\ \hline
\end{tabular}
\end{table}

\subsection{Floating-Point Data Formats}

\subsubsection{Double-Precision Format (IEEE 754-2008 binary64)}
The primary data format supported by the FPU is the 64-bit double-precision binary floating-point format as defined in IEEE 754-2008:

\begin{table}[H]
\centering
\begin{tabular}{|l|c|l|}
\hline
\textbf{Field} & \textbf{Bits} & \textbf{Description} \\ \hline
Sign (S) & 1 & 0 = positive, 1 = negative \\ \hline
Exponent (E) & 11 & Biased exponent, bias = 1023 (0x3FF) \\ \hline
Mantissa / Fraction (M) & 52 & Fractional part of significand \\ \hline
Total & 64 & 64-bit double-precision floating-point word \\ \hline
\end{tabular}
\end{table}

\subsubsection{Single-Precision Format (IEEE 754-2008 binary32)}
The FPU also supports single-precision (32-bit) operations as required by the RISC-V F-extension (RV64F).

\begin{table}[H]
\centering
\begin{tabular}{|l|c|l|}
\hline
\textbf{Field} & \textbf{Bits} & \textbf{Description} \\ \hline
Sign (S) & 1 & 0 = positive, 1 = negative \\ \hline
Exponent (E) & 8 & Biased exponent, bias = 127 (0x7F) \\ \hline
Mantissa / Fraction (M) & 23 & Fractional part \\ \hline
Total & 32 & NaN-boxed to 64-bit in FP registers \\ \hline
\end{tabular}
\end{table}

\subsubsection{Special Values}

\begin{table}[H]
\centering
\begin{tabular}{|l|l|l|l|}
\hline
\textbf{Value} & \textbf{Exponent} & \textbf{Mantissa} & \textbf{Description} \\ \hline
$\pm$ Zero & 0x000 & 0 & Positive and negative zero \\ \hline
Subnormal & 0x000 & Non-zero & Gradual underflow \\ \hline
$\pm$ Infinity & 0x7FF & 0 & Overflow / divide-by-zero \\ \hline
qNaN & 0x7FF & MSB=1 & Quiet NaN \\ \hline
sNaN & 0x7FF & MSB=0 & Signalling NaN \\ \hline
Canonical NaN & 0x7FF & 0x8000000000000 & 0x7FF8000000000000 \\ \hline
\end{tabular}
\end{table}

\subsection{Floating-Point Register File}

\subsubsection{Register Organization}
The FPU provides 32 dedicated 64-bit floating-point registers (f0–f31).

\begin{table}[H]
\centering
\begin{tabular}{|l|l|l|}
\hline
\textbf{Register} & \textbf{ABI Name} & \textbf{Description} \\ \hline
f0–f7 & ft0–ft7 & Temporaries \\ \hline
f8–f9 & fs0–fs1 & Saved \\ \hline
f10–f11 & fa0–fa1 & Arguments/return \\ \hline
f12–f17 & fa2–fa7 & Arguments \\ \hline
f18–f27 & fs2–fs11 & Saved \\ \hline
f28–f31 & ft8–ft11 & Temporaries \\ \hline
\end{tabular}
\end{table}

\subsubsection{Floating-Point Control and Status Register (fcsr)}

\begin{table}[H]
\centering
\begin{tabular}{|l|c|c|l|}
\hline
\textbf{Field} & \textbf{Bits} & \textbf{CSR} & \textbf{Description} \\ \hline
NX & 0 & 0x001 & Inexact \\ \hline
UF & 1 & 0x001 & Underflow \\ \hline
OF & 2 & 0x001 & Overflow \\ \hline
DZ & 3 & 0x001 & Divide by Zero \\ \hline
NV & 4 & 0x001 & Invalid Operation \\ \hline
frm & [7:5] & 0x002 & Rounding Mode \\ \hline
\end{tabular}
\end{table}

\subsection{Rounding Modes}

\begin{table}[H]
\centering
\begin{tabular}{|c|c|l|l|}
\hline
\textbf{Encoding} & \textbf{Mnemonic} & \textbf{IEEE Mode} & \textbf{Description} \\ \hline
000 & RNE & roundTiesToEven & Default \\ \hline
001 & RTZ & roundTowardZero & Truncate \\ \hline
010 & RDN & roundTowardNegative & Floor \\ \hline
011 & RUP & roundTowardPositive & Ceiling \\ \hline
100 & RMM & roundTiesToAway & Away from zero \\ \hline
111 & DYN & -- & Dynamic \\ \hline
\end{tabular}
\end{table}

\subsection{Exception Handling}

\subsubsection{IEEE 754 Exception Flags}

\begin{table}[H]
\centering
\begin{tabular}{|l|c|l|}
\hline
\textbf{Exception} & \textbf{Bit} & \textbf{Condition} \\ \hline
Invalid (NV) & 4 & NaN, invalid ops \\ \hline
Divide by Zero (DZ) & 3 & Non-zero / 0 \\ \hline
Overflow (OF) & 2 & Exceeds range \\ \hline
Underflow (UF) & 1 & Subnormal \\ \hline
Inexact (NX) & 0 & Rounded \\ \hline
\end{tabular}
\end{table}

\subsubsection{Default Results}
\begin{itemize}
\item Invalid: Canonical NaN
\item Divide by Zero: $\pm \infty$
\item Overflow: $\pm \infty$ or max value
\item Underflow: Subnormal or zero
\item Inexact: Rounded result
\end{itemize}

\subsubsection{Exception Flag Encoding Convention}
All five exception flags are packed into a single 5-bit vector by every
arithmetic unit, using a fixed bit order consistent with the \texttt{fflags}
CSR layout:
\[
    \texttt{fflags} = \{\,\text{NV}[4],\ \text{DZ}[3],\ \text{OF}[2],\ \text{UF}[1],\ \text{NX}[0]\,\}
\]
Each functional unit drives this vector combinationally as part of its
result. On instruction commit, the active unit's flag vector is OR-ed
(sticky-accumulated, per IEEE 754 \S7) into the \texttt{fcsr.fflags} field --
flags persist across instructions until explicitly cleared by software.

\subsection{Supported Instructions}

\subsubsection{Arithmetic Instructions}
FADD.D, FSUB.D, FMUL.D, FDIV.D, FSQRT.D, FADD.S, FSUB.S, FMUL.S, FDIV.S, FSQRT.S

\textcolor{blue}{\textbf{\textit{Note: Fused multiply-add instructions
(FMADD.D, FMSUB.D, FNMADD.D, FNMSUB.D) are not yet implemented in the current
RTL. No dedicated fused-multiply-add module exists, and \texttt{fpu\_ctrl.sv}
does not decode these opcodes. They are listed under Future Extensions
(Section~4.2.14).}}}

\subsubsection{Load/Store}
FLD, FSD, FLW, FSW

\textcolor{blue}{\textbf{\textit{Note: these instructions bypass the FPU
arithmetic datapath entirely. They are decoded but routed through the
standard \gls{D-MEM} load/store path shared with integer loads/stores; only
the destination register file (integer vs.\ floating-point) differs.}}}

\subsubsection{Comparison}
FEQ.D, FLT.D, FLE.D, FEQ.S, FLT.S, FLE.S

\subsubsection{Min / Max}
FMIN.D, FMAX.D, FMIN.S, FMAX.S

\textcolor{blue}{\textbf{\textit{Note: implemented as a dedicated unit
separate from the arithmetic core (see Section~4.2.8.2), not as part of the
add/subtract datapath.}}}

\subsubsection{Conversion (Integer $\leftrightarrow$ Floating-Point)}
FCVT.W.D, FCVT.WU.D, FCVT.L.D, FCVT.LU.D, FCVT.D.W, FCVT.D.WU, FCVT.D.L,
FCVT.D.LU, and the corresponding single-precision FCVT.*.S / FCVT.S.* variants

\subsubsection{Conversion (Floating-Point $\leftrightarrow$ Floating-Point)}
FCVT.S.D, FCVT.D.S

\textcolor{blue}{\textbf{\textit{Note: this category was missing from the
previous revision. It performs precision conversion between single and
double formats, including NaN-boxing on widen and Guard/Round/Sticky-based
rounding on narrow.}}}

\subsubsection{Move / Sign-Injection}
FMV.X.D, FMV.D.X, FMV.X.W, FMV.W.X, FSGNJ.D, FSGNJN.D, FSGNJX.D, FSGNJ.S,
FSGNJN.S, FSGNJX.S

\textcolor{blue}{\textbf{\textit{Note: the single-precision move variants
(FMV.X.W, FMV.W.X) and single-precision sign-injection variants were
implemented but omitted from the previous revision's list.}}}

\subsubsection{Classify}
FCLASS.D, FCLASS.S

\subsection{FPU Microarchitecture}

\subsubsection{Upcast-Compute-Downcast Methodology}
Rather than implementing two independent arithmetic datapaths -- one native
to single precision and one native to double precision -- every arithmetic
unit (\texttt{fpu\_addsub}, \texttt{fpu\_mul}, \texttt{fpu\_div},
\texttt{fpu\_sqrt}) implements only the double-precision datapath. A
single-precision operand is upcast to double precision before entering this
shared pipeline; the double-precision result is downcast back to single
precision only when the instruction's format field requests it:

\[
    \text{S operand}
    \xrightarrow{\texttt{s\_to\_d()}}
    \text{D operand}
    \xrightarrow{\text{shared double-precision pipeline}}
    \text{D result}
    \xrightarrow{\texttt{d\_to\_s\_nanbox()}}
    \text{S result (if format = S)}
\]

This is IEEE 754-correct, not merely an implementation shortcut: the
standard defines \texttt{FADD.S} as ``compute the exact mathematical result,
then round once to single precision.'' Double precision carries 53
significand bits versus single's 24 -- sufficient headroom to represent the
exact intermediate result before any rounding occurs. Rounding once, at the
final downcast step, therefore yields the identical, IEEE-correct result
that a native single-precision datapath would produce, with no
double-rounding error.

The \texttt{s\_to\_d}/\texttt{d\_to\_s\_nanbox} conversion functions are
implemented locally within each of the four arithmetic modules rather than
factored into a shared module -- a deliberate tradeoff of minor code
duplication in exchange for fully self-contained, independently verifiable
units.

\subsubsection{Pipeline Integration}
Fetch → Decode → FP Execute → Writeback

\subsubsection{Execution Latencies}
\textcolor{red}{\textbf{Correction from previous revision:}} the FPU does
not currently implement multi-cycle execution for any operation. All
eleven functional units (\texttt{fpu\_addsub}, \texttt{fpu\_mul},
\texttt{fpu\_div}, \texttt{fpu\_sqrt}, \texttt{fpu\_cmp}, \texttt{fpu\_cvt},
\texttt{fpu\_cvt\_fmt}, \texttt{fpu\_sign}, \texttt{fpu\_minmax},
\texttt{fpu\_mv}, \texttt{fpu\_classify}) are purely combinational and
execute within a single clock cycle, alongside the integer \gls{alu}, inside
the \texttt{EXEC} state of the \texttt{top\_npl} \gls{fsm}.

\begin{table}[H]
\centering
\caption{FPU Unit Latency (Current Implementation)}
\begin{tabular}{|l|l|l|}
\hline
\textbf{Unit} & \textbf{Latency} & \textbf{Note} \\ \hline
fpu\_addsub, fpu\_mul, fpu\_cmp, fpu\_cvt, & 0 cycles & Fully combinational \\
fpu\_cvt\_fmt, fpu\_sign, fpu\_minmax, & & \\
fpu\_mv, fpu\_classify & & \\ \hline
fpu\_div & 0 cycles & Combinational; exposes
\texttt{clk}/\texttt{valid\_in}/\texttt{valid\_out}/\texttt{ready\_out}
ports for a future pipelined revision, but \texttt{ready\_out} is hardwired
high and \texttt{valid\_out = valid\_in \& en} -- no internal state or
iteration exists today. \\ \hline
fpu\_sqrt & 0 cycles & Non-restoring bit-by-bit algorithm, unrolled as a
combinational loop (53 stages) rather than iterated over clock cycles. \\
\hline
\end{tabular}
\end{table}

\textcolor{blue}{\textbf{\textit{Synthesis caveat: the combinational divider
and square-root units produce a long critical path and are expected to
require timing closure work (or conversion to a multi-cycle/pipelined
implementation using the existing handshake ports) for any target above
simulation. This is a documented, known limitation of the current
implementation, not an oversight.}}}

A future pipelined revision targeting non-zero latencies (e.g.\ FADD: 3,
FMUL: 4, FDIV: 20--24, FSQRT: 25--30 cycles) is tracked under Future
Extensions (Section~4.2.14) and is not part of the current implementation.

\subsubsection{Functional Unit Map}
The FPU is composed of eleven independently instantiated submodules,
dispatched from a single combinational control block (\texttt{fpu\_ctrl}).
All units execute in parallel every cycle; a result multiplexer, driven by
the decoded opcode, selects the active unit's output.

\begin{table}[H]
\centering
\caption{FPU Functional Unit Map}
\begin{tabular}{|l|l|l|}
\hline
\textbf{Module} & \textbf{Instructions Handled} & \textbf{Output Destination} \\ \hline
fpu\_addsub  & FADD, FSUB                          & FP regfile \\ \hline
fpu\_mul     & FMUL                                & FP regfile \\ \hline
fpu\_div     & FDIV                                & FP regfile \\ \hline
fpu\_sqrt    & FSQRT                               & FP regfile \\ \hline
fpu\_cmp     & FEQ, FLT, FLE                       & Integer regfile \\ \hline
fpu\_cvt     & FCVT.* (int $\leftrightarrow$ float) & Integer or FP regfile \\ \hline
fpu\_cvt\_fmt& FCVT.S.D, FCVT.D.S                  & FP regfile \\ \hline
fpu\_sign    & FSGNJ, FSGNJN, FSGNJX               & FP regfile \\ \hline
fpu\_minmax  & FMIN, FMAX                          & FP regfile \\ \hline
fpu\_mv      & FMV.X.D, FMV.D.X, FMV.X.W, FMV.W.X  & Integer or FP regfile \\ \hline
fpu\_classify& FCLASS                              & Integer regfile \\ \hline
\end{tabular}
\end{table}

The destination register file is selected by the \texttt{to\_int} signal,
described in Section~4.2.9.

\subsubsection{FP Register File}
32 × 64-bit, 2 read ports, 1 write port

\subsubsection{Operand Forwarding}
Supported for single-cycle ops

\subsubsection{NaN Handling}
IEEE 754 compliant

\subsubsection{Subnormal Handling}
Fully supported (no flush-to-zero)

\subsection{Hardware Interfaces}
\textcolor{red}{\textbf{Correction from previous revision:}} the signal
names listed previously (\texttt{fpu\_op}, \texttt{fpu\_rs1},
\texttt{fpu\_rs2}, \texttt{fpu\_ready}) do not correspond to the actual
\texttt{fpu\_top} port list. The real interface is:

\begin{table}[H]
\centering
\caption{fpu\_top Port List}
\begin{tabular}{|l|l|}
\hline
\textbf{Signal} & \textbf{Description} \\ \hline
\texttt{inst} & Decoded instruction enum, shared with the integer decode stage \\ \hline
\texttt{fp\_funct5} & Sub-operation select field (extracted from the opcode) \\ \hline
\texttt{fp\_rm} & Rounding mode (resolved from \texttt{fcsr.frm} if dynamic) \\ \hline
\texttt{fp\_cvt\_toint} & Selects integer-destination conversion direction \\ \hline
\texttt{fp\_cvt\_word} & Selects 32-bit vs.\ 64-bit integer conversion width \\ \hline
\texttt{fp\_cvt\_signed} & Selects signed vs.\ unsigned integer conversion \\ \hline
\texttt{fp\_fmt} & Selects single vs.\ double precision format \\ \hline
\texttt{fp\_sgn\_op} & Selects FSGNJ / FSGNJN / FSGNJX variant \\ \hline
\texttt{fp\_min\_sel} & Selects FMIN vs.\ FMAX \\ \hline
\texttt{operand\_a, operand\_b} & 64-bit floating-point operands \\ \hline
\texttt{int\_operand} & Integer operand (for int$\to$float conversion) \\ \hline
\texttt{result} & 64-bit result (floating-point or integer, per \texttt{to\_int}) \\ \hline
\texttt{to\_int} & Selects integer regfile (1) vs.\ FP regfile (0) as write-back destination \\ \hline
\texttt{fflags} & 5-bit exception flag vector, see Section~4.2.6 \\ \hline
\texttt{div\_valid\_in, div\_ready,} & Handshake stub ports for the divider's future pipelined \\
\texttt{div\_valid\_out} & revision; currently always-ready/combinational (see Section~4.2.8) \\ \hline
\end{tabular}
\end{table}
\subsection{Floating-Point CSR Registers}

\begin{table}[H]
\centering
\begin{tabular}{|c|l|c|c|c|l|}
\hline
\textbf{Addr} & \textbf{Name} & \textbf{Width} & \textbf{Access} & \textbf{Reset} & \textbf{Description} \\ \hline
0x001 & fflags & 32 & RW & 0 & Exception flags \\ \hline
0x002 & frm & 32 & RW & 0 & Rounding mode \\ \hline
0x003 & fcsr & 32 & RW & 0 & Combined \\ \hline
\end{tabular}
\end{table}

\subsection{Address Map and Integration}
FPU is internal, not memory-mapped.

\subsection{Power Considerations}
Always powered; clock gating planned.

\subsection{Integration with top\_npl}
The FPU is integrated as a second execution unit running in parallel with
the integer \gls{alu}, sharing the same decode stage and committing in the
same cycle.

\begin{itemize}
    \item \textbf{Shared decode:} the decoded instruction feeds both
          \texttt{alu\_top} and \texttt{fpu\_top} in the same cycle; no
          separate FP decode stage exists.
    \item \textbf{Parallel register file reads:} a dedicated 32$\times$64-bit
          FP register file is read every cycle alongside the integer
          register file, both asynchronously, so floating-point operands
          are available with no additional read latency.
    \item \textbf{Single-cycle execution:} \texttt{fpu\_top} evaluates
          combinationally within the \texttt{EXEC} state of the
          \texttt{top\_npl} \gls{fsm} -- the same cycle used by the integer
          \gls{alu}.
    \item \textbf{Write-back routing:} the \texttt{to\_int} signal selects
          the destination register file via two mutually exclusive
          write-enable equations:
\end{itemize}

\begin{verbatim}
int_we = commit && !trap && reg_write &&
         (!is_fp || (is_fp && to_int));

fp_we  = commit && !trap && reg_write &&
         (fp_load || (is_fp && !fp_load && !to_int));
\end{verbatim}

\begin{itemize}
    \item \textbf{Load/store bypass:} FLD/FSD are decoded to
          \texttt{FPU\_SEL\_PASS} in \texttt{fpu\_ctrl} and do not enter the
          FPU arithmetic datapath; they reuse the same \gls{D-MEM}
          load/store path as integer loads/stores, differing only in
          destination register file.
    \item \textbf{Exception flag accumulation:} on every committed
          floating-point instruction, the active unit's \texttt{fflags}
          output is OR-ed into \texttt{fcsr.fflags} (sticky accumulation,
          per Section~4.2.6).
\end{itemize}

\subsection{Verification and Test Strategy}

\subsubsection{Unit-Level Simulation}
SystemVerilog testbenches

\subsubsection{IEEE Compliance Testing}
Berkeley TestFloat

\subsubsection{Integration Testing}
RV64I toolchain tests

\subsubsection{Regression Testing}
CTest framework

\subsection{Future Extensions}
\begin{itemize}
\item Fused multiply-add instructions (FMADD.D, FMSUB.D, FNMADD.D, FNMSUB.D) -- not yet implemented
\item Pipelined/multi-cycle FDIV and FSQRT using the existing handshake ports, targeting the latencies listed in Section~4.2.8
\item Flush-to-Zero
\item Quad precision (Q-extension)
\item Clock gating
\item Hardware exception trapping
\item Transcendental functions
\end{itemize}
